\section{Related Work}
\label{sec:related-work}

% DOCX paragraph 021.
\paragraph{Planning and tool use.}
Existing methods organize tool calls from different perspectives. ReAct interleaves reasoning, actions, and observation feedback~\cite{yao2023react}. ToolPlanner constructs invocation workflows using path planning and feedback~\cite{wu2024toolplanner}, whereas Tool-Planner groups tools by function to support task planning~\cite{liu2025toolplanner}. LLM+P separates language processing from classical planning search~\cite{liu2023llmp}, and Generalised BDI Planning integrates planning with agent goal management and action selection~\cite{meneguzzi2025bdi}. AgentGS likewise combines language models with structured planning. It uses capabilities and dependencies declared in registered Skills to organize required operations and fill missing steps before selecting concrete implementations.

% DOCX paragraph 022. Source references 013--015 are incorporated below.
\paragraph{Tool relations and capability composition.}
Representations of tools and their dependencies provide structural support for multistep invocation. ToolNet organizes tools into a directed graph for tool navigation~\cite{liu2024toolnet}, and ControlLLM searches a graph to compose tools~\cite{liu2024controlllm}. In-N-Out studies parameter-level API graphs, emphasizing the relationships involved in passing arguments between calls~\cite{lee2025innout}. HyperAgent constructs a hypergraph from schemas for tool inputs and outputs and searches for supporting tools when requirements remain unsatisfied~\cite{zhai2026hyperagent}. AgentGS instead organizes strategy knowledge with shared Capabilities as nodes and Skills as hyperedges, and links concrete implementations through mappings between Capabilities and Tools. This representation preserves a capability's dependencies and usage constraints within each Skill. It supports reuse of local operations across strategies and retrieval of connecting branches along dependencies carrying their Skill provenance, thereby completing missing capabilities in the current plan.

% DOCX paragraph 023.
\paragraph{Geoscientific evaluation.}
GeoBench, introduced with K2, evaluates knowledge understanding through examination-style geoscientific questions~\cite{deng2023k2}. EarthSE includes professional questions and open-ended scientific dialogue to assess geoscientific exploration~\cite{xu2025earthse}. These evaluations may require multistep reasoning, but do not necessarily require multiple actual tool invocations. SciToolAgent, GeoAgent, and Earth-Agent execute scientific data processing and analysis through tool use~\cite{ding2025scitoolagent,chen2024geoagent,feng2026earthagent}.

% DOCX paragraph 024.
\paragraph{Multitool orchestration benchmarks.}
GeoBenchX requires agents to use geospatial tools for multistage tasks and evaluates results against reference solutions~\cite{krechetova2025geobenchx}. Earth-Bench distinguishes autonomous planning from instruction following and records tool inputs and outputs~\cite{feng2026earthagent}. GeoAgentBench examines invocation order, parameters, and final maps in a GIS sandbox~\cite{yu2026geoagentbench}. GeoNatureAgent Benchmark evaluates multitool analysis and error recovery using real APIs~\cite{diazireland2026geonatureagent}. Existing benchmarks thus cover multistep tool use and some dependency verification. We further investigate the composition of scientific operations across geoscientific themes. Alongside tool use and final answers, we examine agreement with reference invocation sequences, actual use of intermediate artifacts, and completion of necessary operations. GeoToolEval emphasizes scientific-operation composition and execution-process analysis. A detailed comparison with existing benchmarks is reserved for Supplementary Material~I, \drafttodo{benchmark comparison material to be supplied}.
